\bookchapter{الفصل 48: التعرف على أندرويد — من نظام التشغيل إلى منظومة التطوير}

\chapterintro{مبروك عليك هذا القرار — أن تتحول من “شخص يلعب بالهاتف” إلى “شخص يصنع تطبيقات الهاتف”. وقبل أن نشمر عن الأكمام، لننشر الخريطة أولا: من أين جاء أندرويد؟ لماذا يكتب له المطورون في العالم كله التطبيقات؟ وماذا يحدث بالضبط بين لحظة ضربك المفاتيح ووصول التطبيق إلى هواتف المستخدمين؟}

\booksection{أولاً: قصة أندرويد — من “الروبوت الصغير” إلى “أكبر أنظمة العالم”}

تبدأ قصة أندرويد (Android) عام 2003 — حين أسس مهندس اسمه Andy Rubin شركة Android بهدف بسيط: بناء نظام تشغيل ذكي “يفهم الكاميرا أكثر”. ثم استحوذت عليه جوجل عام 2005، وفي 2007 شكلت مع عشرات المصنعين تحالف الهواتف المفتوحة (OHA)، وفي سبتمبر 2008 طرح أول هاتف أندرويد رسميا وهو HTC Dream (G1).

قصة ما بعدها يعرفها الجميع: بحلول 2026 أصبح أندرويد نظام التشغيل المحمول الأوسع حصة سوقية في العالم، وعدد الأجهزة النشطة بمليارات. هاتفك وعلبة التلفاز وشاشة السيارة وساعتك الذكية — كلها على الأرجح تعمل بأندرويد. ماذا يعني هذا؟ يعني أن كل سطر تكتبه من الكود له فرصة أن يستخدمه أكبر عدد من الناس في العالم.

\begin{tcbinfo}{نصيحة: الجينات مفتوحة المصدر في أندرويد}
الكود الأساسي لأندرويد (AOSP، مشروع أندرويد مفتوح المصدر) مفتوح المصدر — وأي مصنع يستطيع أن يأخذه ويعدله ويخصصه. لذلك لدى Xiaomi واجهة MIUI (وتسمى الآن HyperOS)، ولدى Huawei طبقة توافق HarmonyOS، ولدى Samsung واجهة One UI. وهذا خير للمطورين: أنت تتعلم “تطوير أندرويد” لا “تطوير علامة بعينها” — مهارة واحدة تغطي كل الأجهزة.

\end{tcbinfo}

وتطور أرقام الإصدارات ممتع أيضا: بدءا من 1.5 باسم Cupcake (الكب كيك)، سمي كل إصدار كبير باسم حلوى — Donut وÉclair وFroyo وGingerbread... إلى أن جاء التحول إلى الأرقام الخالصة مع أندرويد 10 عام 2019. وحتى أغسطس 2026، الإصدار الرسمي الأحدث هو Android 17، أما “الحد الأدنى للدعم” الذي يهمك فعلا عند التطوير فتضعه التطبيقات السائدة عادة بين Android 8 (API 26) وAndroid 10 (API 29) — كلما دعمت أجهزة السنوات الماضية اتسعت قاعدة مستخدميك.

\begin{center}
\begin{tabularx}{\linewidth}{|X|X|X|}
\hline
\textbf{الإصدار} & \textbf{الاسم الرمزي/السنة} & \textbf{ما يهم المطور} \tabularnewline
\hline
Android 8 (API 26) & Oreo · 2017 & قنوات الإشعارات وقيود العمل في الخلفية — الحد الأدنى للعديد من الأجهزة القديمة \tabularnewline
Android 10 (API 29) & 2019 & ظهور التخزين المقسم وتغيير جذري في قواعد الوصول إلى الملفات \tabularnewline
Android 12 (API 31) & 2021 & ألوان ديناميكية بأسلوب Material You وترقية كبيرة لمظهر الواجهة \tabularnewline
Android 13 (API 33) & 2022 & إذن الإشعارات وقت التشغيل، ونافذة تسألك "هل تسمح بإرسال الإشعارات" \tabularnewline
Android 14-15 (API 34-35) & 2023-2024 & الرجوع التنبؤي، وصفحات 16KB، وتشديد إضافي للخصوصية \tabularnewline
Android 16-17 (API 36-37) & 2025-2026 & إيقاع أسرع للإصدارات الصغيرة، ودمج قدرات الذكاء الاصطناعي (مثل توسعات Gemini) \tabularnewline
\hline
\end{tabularx}
\end{center}

\begin{center}\small\color{gray}تاريخ موجز لإصدارات أندرويد (حتى أغسطس 2026، النقاط التي يركز عليها المطورون)\end{center}

\booksection{ثانياً: بأي لغة نكتب — Java أم Kotlin؟}

ظهرت في تاريخ تطوير أندرويد لغتان رئيسيتان. Java هي “العمود الفقري المؤسس” — فقد استخدمت منذ ولادة أندرويد؛ وKotlin هي “المواطنة من الدرجة الأولى” التي اختارتها جوجل عام 2017، ومنذ 2019 أعلنت رسميا مبدأ “Kotlin أولا” (Kotlin-first)، وبحلول 2026 صارت المشاريع الجديدة كلها تقريبا بـ Kotlin، وترحل المشاريع القديمة معظمها إليها تباعا.

وبصراحة: إذا كنت مبتدئا تاما فتعلّم Kotlin مباشرة ولا ترجع إلى Java. ثلاثة أسباب: أولا كود Kotlin أقل من Java بنسبة 30\%-40\% (نفس الوظيفة لكنه أقصر وأوضح)؛ ثانيا يقيك بطبيعته من “المؤشر الفارغ” (Null Pointer) كابوس المبتدئين؛ ثالثا كل الدروس الحديثة والوثائق الرسمية وأدوات البرمجة بالذكاء الاصطناعي (أتذكر كتابنا في البرمجة بالذكاء الاصطناعي؟) تعتمد Kotlin مرجعا. أما Java فلن تصادفها إلا في صيانة المشاريع القديمة وقراءة الكود القديم.

أسلوب Java (النسخة المطولة)

\begin{lstlisting}
public class MainActivity extends Activity \{
@Override
protected void onCreate(Bundle b) \{
super.onCreate(b);
setContentView(R.layout.activity);
\}
\}
\end{lstlisting}

أسلوب Kotlin (النسخة المقتضبة)

\begin{lstlisting}
class MainActivity : ComponentActivity() \{
override fun onCreate(b: Bundle?) \{
super.onCreate(b)
setContent \{ App() \}
\}
\}
\end{lstlisting}

نفس العملية تماما، لكن Kotlin يستغني عن الكثير من الأقواس المعقوفة والكود النمطي. سنتعلمه رسميا في الفصل 50، والآن يكفيك أن تحفظ: “اللغة العامة” لتطوير أندرويد اسمها Kotlin.

\booksection{ثالثاً: كيف يولد التطبيق (App)؟}

يظن كثيرون أن “تطوير التطبيقات” مجرد فتح الحاسوب والكتابة بلا توقف. في الواقع يحتاج التطبيق من الفكرة إلى هواتف المستخدمين إلى خط إنتاج تقريبي كالتالي:

فكرة

التصميم (الواجهة + الوظائف)

كتابة الكود

الترجمة والتغليف (APK/AAB)

الاختبار (محاكي + جهاز حقيقي)

النشر (المتجر/حزمة التثبيت)

\begin{center}\small\color{gray}الشكل 1: الدورة الكاملة لتطبيق أندرويد من الفكرة إلى النشر في المتجر\end{center}

في هذا المسار مفهومان عليك التعرف عليهما من الآن:

\begin{itemize}
  \item APK: صيغة حزمة تثبيت تطبيقات أندرويد (.apk)، تنقر عليها فيثبت في الهاتف، تماما مثل ملف .exe في Windows. استخدمها لاختباراتك الشخصية ولمشاركة التطبيق مع أصدقائك للتجربة.
  \item AAB: صيغة Android App Bundle (.aab)، وهي "صيغة التغليف المتقدمة" المستخدمة للنشر على Google Play. يقوم المتجر بقص الحزمة ديناميكيا حسب طراز هاتف كل مستخدم، فيكون ما ينزله كل شخص أصغر حجما. لا يمكنك تثبيتها بنفسك، وهي تسلم إلى المتجر فقط.
\end{itemize}

وسوف تسمع باستمرار مصطلحا آخر: SDK (مجموعة تطوير البرمجيات). وهي “عدة التطوير” الكاملة التي توفرها أندرويد رسميا — تشمل مكتبات API لكل الإصدارات وأدوات التصحيح وصور المحاكيات وغيرها. الكلاسات مثل Activity وButton التي تستدعيها في كودك تأتي كلها من الـ SDK.

\booksection{رابعاً: المشهد الكامل لحزمة تقنيات أندرويد — “قائمة أسلحتك”}

قبل البدء بالتعلم الرسمي، تعرف على الإعداد القياسي لمشاريع أندرويد السائدة في 2026 — لا يهم إن لم تفهم هذه القائمة الآن، فعندما تعود إليها بعد إنهاء الكتاب كلها ستنفتح أمامك تماما:

\begin{center}
\begin{tabularx}{\linewidth}{|X|X|X|}
\hline
\textbf{الطبقة} & \textbf{الخيار السائد (2026)} & \textbf{لأي غرض يستخدم} \tabularnewline
\hline
اللغة & Kotlin 2.x & لغة كتابة المنطق \tabularnewline
الواجهة & Jetpack Compose + Material 3 & رسم الواجهة بالكود (الخط الرئيسي لهذا الكتاب) \tabularnewline
البنية المعمارية & MVVM + ViewModel & فصل الواجهة عن المنطق حتى لا يتشابك الكود \tabularnewline
تخزين البيانات & Room (قاعدة بيانات) + DataStore (الإعدادات) & حفظ البيانات في الهاتف \tabularnewline
الشبكات & Retrofit + OkHttp + Coroutines في Kotlin & التحدث مع الخادم \tabularnewline
أدوات البناء & Gradle + Kotlin DSL & تحويل الكود إلى حزمة تثبيت \tabularnewline
بيئة التطوير & Android Studio & بيئة التطوير المتكاملة "الشاملة" الرسمية من جوجل \tabularnewline
\hline
\end{tabularx}
\end{center}

\begin{center}\small\color{gray}نظرة سريعة على حزمة تقنيات أندرويد السائدة في 2026\end{center}

لا تدع هذه السلسلة من الكلمات الأجنبية تخيفك. إنها أشبه بأدوات المطبخ — كل ما تحتاج معرفته الآن هو أن “الطبخ يجري في المطبخ”، وكيفية استخدام كل أداة سنفككها لك واحدة تلو الأخرى في الفصول القادمة. والخبر السار: عند إنشاء مشروع جديد يهيئ Android Studio هذه المجموعة كلها تلقائيا، وأنت مكلف بالطبخ فقط.

\booksection{خامساً: أندرويد مقابل iOS — أول درس “اختيار المعسكر” للمطور}

قد تكون تحمل الآن جهاز أندرويد وجهاز آبل معا. بصفتك مطورا، من المفيد التعرف مسبقا على الفروق بين التجربتين — ليس انحيازا لسؤال “أيهما أفضل” بل حسابا واقعيا لسؤال “أي الجهتين أنسب لبدايتك”:

\begin{center}
\begin{tabularx}{\linewidth}{|X|X|X|}
\hline
\textbf{البعد} & \textbf{تطوير أندرويد} & \textbf{تطوير iOS} \tabularnewline
\hline
حاسوب التطوير & Windows / Mac / Linux — أي منها يصلح & يتطلب جهاز Mac حصرًا (Xcode متاح عليه فقط) \tabularnewline
أدوات التطوير & Android Studio (مجاني) & Xcode (مجاني) + رسم سنوي للنشر قدره 99 دولارًا \tabularnewline
مراجعة النشر على المتجر & متساهلة، من بضع ساعات إلى بضعة أيام & صارمة، ومن 1 إلى 7 أيام هو المعتاد \tabularnewline
التثبيت والتجربة على الأجهزة & ملف APK يمكن نشره وتثبيته بحرية تامة & التوزيع محدود، والاختبار يتطلب التوقيع أو TestFlight \tabularnewline
بيئة تحقيق الإيرادات & الأغلبية للإعلانات والقنوات المحلية & عادات الدفع داخل App Store أفضل (المستخدمون أكثر استعدادًا للشراء داخل التطبيق) \tabularnewline
كلفة التعلم & مفاهيم أكثر (المكونات الأربعة وغيرها) & نظام مغلق لكنه موحد وبتوثيق ممتاز \tabularnewline
\hline
\end{tabularx}
\end{center}

\begin{center}\small\color{gray}مقارنة بين مساري تطوير Android و iOS (أغسطس 2026)\end{center}

لنتحدث بصراحة: إذا لم يكن حاسوبك من نوع Mac، فلا يوجد أصلًا سؤال للتفكير فيه — Android هو الخيار الوحيد. أما إذا كان لديك Mac، فيمكنك قراءة الجزأين الثامن والتاسع من هذا الكتاب بالتوازي، فالأفكار البرمجية واحدة، وبتعلم أحد المسارين سيصبح تعلم الآخر أسرع بكثير.

\begin{tcbpractice}{لنتجرب: انظر إلى نظام Android في هاتفك}
افتح في هاتفك [الإعدادات] > [حول الهاتف] (أو «حول هذا الجهاز»)، وابحث عن السطر الخاص بـ«إصدار Android» — وانقر عليه عدة نقرات متتابعة لتفعيل لعبة المفاجآت المخفية في النظام. دوّن رقم الإصدار لديك وقارنه بالجدول أعلاه لترى إلى أي جيل ينتمي. هاتفك هو أول «جهاز اختبار للتطوير» لديك.

\end{tcbpractice}

\begin{tcbsummary}{خلاصة الفصل}
\begin{itemize}
  \item وُلد Android عام 2008، وفي 2026 أصبح النظام المحمول الأكبر حصةً سوقية عالميًا بمليارات الأجهزة النشطة
  \item أحدث إصدار رسمي هو Android 17؛ والتطبيقات الرئيسية تدعم كحد أدنى Android 8-10 لتكون قاعدة المستخدمين واسعة بما يكفي
  \item استخدم Kotlin في كل مشروع جديد: كود أقل بنسبة 30\%+، وحماية من المؤشرات الفارغة، وهو الاختيار الرسمي المعتمد
  \item ملف APK تثبته بنفسك، وصيغة AAB للمتجر؛ وSDK هي حزمة أدوات التطوير الرسمية
  \item حزمة التقنيات القياسية لعام 2026: Kotlin + Compose + MVVM + Room + Retrofit + Gradle + Android Studio
\end{itemize}

\end{tcbsummary}
